% GENERATED FILE. Edit canonical lessons, then run npm run generate:books.
% canonical-source-hash: fnv1a64:4a32620a8a16e49a
% canonical-lessons: TA-C224-minavar, TA-C224-kuyavar, TA-C224-mutitiruttupavar, TA-C224-kottanar, TA-C224-vinnani, TA-R224-first-pass-words-from-chapters-221-222, TA-R224-second-pass-words-from-chapters-223-224

\chapter{Fisherman, Potter, Barber, Mason, Scientist}
\label{ch:ta-fisherman-potter-barber-mason-scientist}
\hlchaptermodality{224}

\begin{chapteropening}
\textbf{By the end of this chapter:} \emph{I can say a fisherman, a potter, a barber, a mason and a scientist.}

Complete the last lesson of chapter 224: I can say a fisherman, a potter, a barber, a mason and a scientist.
\end{chapteropening}

\section[\texorpdfstring{\ta{மீனவர்} (mīṉavar)}{மீனவர் (mīṉavar)}]{\ta{மீனவர்} (mīṉavar) — a fisherman}

\label{lesson:TA-C224-minavar}



\begin{quote}
Before the new one: say the Tamil for a lesson; a school subject, then the Tamil for homework.
\end{quote}

\subsection*{You'll want to know: \ta{மீனவர்}}
\textbf{\ta{மீனவர்}} — \emph{mīṉavar} — \textquotedblleft{}a fisherman\textquotedblright{}.

\begin{grammarlens}[title={What you've built}]
One more word for this chapter.
\end{grammarlens}

\subsection*{Your turn}
\begin{itemize}
\raggedright
  \item \emph{Say it:} \emph{mīṉavar}
  \item \emph{Say it:} \emph{mīṉavar}, once more
  \item \emph{Say it:} say \emph{mīṉavar}
  \item \emph{Recall:} say the Tamil for a lesson; a school subject, then the Tamil for homework, then say \emph{mīṉavar} again
\end{itemize}

\begin{tcolorbox}[breakable,colback=teal!4,colframe=teal!35!black,title={Before you move on}]
What does \ta{மீனவர்} mean? (\textquotedblleft{}A fisherman\textquotedblright{}.) Say it once more.
\end{tcolorbox}
\section[\texorpdfstring{\ta{குயவர்} (kuyavar)}{குயவர் (kuyavar)}]{\ta{குயவர்} (kuyavar) — a potter}

\label{lesson:TA-C224-kuyavar}



\begin{quote}
Before the new one: say the Tamil for homework, then the Tamil for a fisherman.
\end{quote}

\subsection*{You'll want to know: \ta{குயவர்}}
\textbf{\ta{குயவர்}} — \emph{kuyavar} — \textquotedblleft{}a potter\textquotedblright{}.

\begin{grammarlens}[title={What you've built}]
One more word for this chapter.
\end{grammarlens}

\subsection*{Your turn}
\begin{itemize}
\raggedright
  \item \emph{Say it:} \emph{kuyavar}
  \item \emph{Say it:} \emph{kuyavar}, once more
  \item \emph{Say it:} say \emph{kuyavar}
  \item \emph{Recall:} say the Tamil for homework, then the Tamil for a fisherman, then say \emph{kuyavar} again
\end{itemize}

\begin{tcolorbox}[breakable,colback=teal!4,colframe=teal!35!black,title={Before you move on}]
What does \ta{குயவர்} mean? (\textquotedblleft{}A potter\textquotedblright{}.) Say it once more.
\end{tcolorbox}
\section[\texorpdfstring{\ta{முடிதிருத்துபவர்} (muṭitiruttupavar)}{முடிதிருத்துபவர் (muṭitiruttupavar)}]{\ta{முடிதிருத்துபவர்} (muṭitiruttupavar) — a barber}

\label{lesson:TA-C224-mutitiruttupavar}



\begin{quote}
Before the new one: say the Tamil for a fisherman, then the Tamil for a potter.
\end{quote}

\subsection*{You'll want to know: \ta{முடிதிருத்துபவர்}}
\textbf{\ta{முடிதிருத்துபவர்}} — \emph{muṭitiruttupavar} — \textquotedblleft{}a barber\textquotedblright{}.

\begin{grammarlens}[title={What you've built}]
One more word for this chapter.
\end{grammarlens}

\subsection*{Your turn}
\begin{itemize}
\raggedright
  \item \emph{Say it:} \emph{muṭitiruttupavar}
  \item \emph{Say it:} \emph{muṭitiruttupavar}, once more
  \item \emph{Say it:} say \emph{muṭitiruttupavar}
  \item \emph{Recall:} say the Tamil for a fisherman, then the Tamil for a potter, then say \emph{muṭitiruttupavar} again
\end{itemize}

\begin{tcolorbox}[breakable,colback=teal!4,colframe=teal!35!black,title={Before you move on}]
What does \ta{முடிதிருத்துபவர்} mean? (\textquotedblleft{}A barber\textquotedblright{}.) Say it once more.
\end{tcolorbox}
\section[\texorpdfstring{\ta{கொத்தனார்} (kottaṉār)}{கொத்தனார் (kottaṉār)}]{\ta{கொத்தனார்} (kottaṉār) — a mason, a bricklayer}

\label{lesson:TA-C224-kottanar}



\begin{quote}
Before the new one: say the Tamil for a potter, then the Tamil for a barber.
\end{quote}

\subsection*{You'll want to know: \ta{கொத்தனார்}}
\textbf{\ta{கொத்தனார்}} — \emph{kottaṉār} — \textquotedblleft{}a mason, a bricklayer\textquotedblright{}.

\begin{grammarlens}[title={What you've built}]
One more word for this chapter.
\end{grammarlens}

\subsection*{Your turn}
\begin{itemize}
\raggedright
  \item \emph{Say it:} \emph{kottaṉār}
  \item \emph{Say it:} \emph{kottaṉār}, once more
  \item \emph{Say it:} say \emph{kottaṉār}
  \item \emph{Recall:} say the Tamil for a potter, then the Tamil for a barber, then say \emph{kottaṉār} again
\end{itemize}

\begin{tcolorbox}[breakable,colback=teal!4,colframe=teal!35!black,title={Before you move on}]
What does \ta{கொத்தனார்} mean? (\textquotedblleft{}A mason, a bricklayer\textquotedblright{}.) Say it once more.
\end{tcolorbox}
\section[\texorpdfstring{\ta{விஞ்ஞானி} (viññāṉi)}{விஞ்ஞானி (viññāṉi)}]{\ta{விஞ்ஞானி} (viññāṉi) — a scientist}

\label{lesson:TA-C224-vinnani}



\begin{quote}
Before the new one: say the Tamil for a barber, then the Tamil for a mason.
\end{quote}

\subsection*{You'll want to know: \ta{விஞ்ஞானி}}
\textbf{\ta{விஞ்ஞானி}} — \emph{viññāṉi} — \textquotedblleft{}a scientist\textquotedblright{}.

\begin{grammarlens}[title={What you've built}]
One more word for this chapter.
\end{grammarlens}

\subsection*{Your turn}
\begin{itemize}
\raggedright
  \item \emph{Say it:} \emph{viññāṉi}
  \item \emph{Say it:} \emph{viññāṉi}, once more
  \item \emph{Say it:} say \emph{viññāṉi}
  \item \emph{Recall:} say the Tamil for a barber, then the Tamil for a mason, then say \emph{viññāṉi} again
\end{itemize}

\begin{tcolorbox}[breakable,colback=teal!4,colframe=teal!35!black,title={Before you move on}]
What does \ta{விஞ்ஞானி} mean? (\textquotedblleft{}A scientist\textquotedblright{}.) Say it once more.
\end{tcolorbox}
\section[(dialogue)]{Words from chapters 221-222 — a first pass}

\label{lesson:TA-R224-first-pass-words-from-chapters-221-222}



\begin{quote}
Say the words for a mason and a scientist.
\end{quote}

\subsection*{Your turn}
Say these aloud:

\begin{itemize}
\raggedright
  \item a helmet, a wheel, a blackboard, an eraser, a uniform
  \item a discount, a rupee, a coin, an expense, a salary
\end{itemize}

\begin{tcolorbox}[breakable,colback=teal!4,colframe=teal!35!black,title={Before you move on}]
A helmet? (\textbf{\ta{தலைக்கவசம்}}, \emph{talaikkavacam}.) A wheel? (\textbf{\ta{சக்கரம்}}, \emph{cakkaram}.) A blackboard? (\textbf{\ta{கரும்பலகை}}, \emph{karumpalakai}.) An eraser? (\textbf{\ta{அழிப்பான்}}, \emph{aḻippāṉ}.) A uniform? (\textbf{\ta{சீருடை}}, \emph{cīruṭai}.) A discount? (\textbf{\ta{தள்ளுபடி}}, \emph{taḷḷupaṭi}.) A rupee? (\textbf{\ta{ரூபாய்}}, \emph{rūpāy}.) A coin? (\textbf{\ta{நாணயம்}}, \emph{nāṇayam}.) An expense? (\textbf{\ta{செலவு}}, \emph{celavu}.) A salary? (\textbf{\ta{சம்பளம்}}, \emph{campaḷam}.)
\end{tcolorbox}
\section[(dialogue)]{Words from chapters 223-224 — a second pass}

\label{lesson:TA-R224-second-pass-words-from-chapters-223-224}



\begin{quote}
Say the words for an interview and a company.
\end{quote}

\subsection*{Your turn}
Say these aloud:

\begin{itemize}
\raggedright
  \item an interview, a company, a factory, a lesson; a school subject, homework
  \item a fisherman, a potter, a barber, a mason, a scientist
\end{itemize}

\begin{tcolorbox}[breakable,colback=teal!4,colframe=teal!35!black,title={Before you move on}]
An interview? (\textbf{\ta{நேர்காணல்}}, \emph{nērkāṇal}.) A company? (\textbf{\ta{நிறுவனம்}}, \emph{niṟuvaṉam}.) A factory? (\textbf{\ta{தொழிற்சாலை}}, \emph{toḻiṟcālai}.) A lesson; a school subject? (\textbf{\ta{பாடம்}}, \emph{pāṭam}.) Homework? (\textbf{\ta{வீட்டுப்பாடம்}}, \emph{vīṭṭuppāṭam}.) A fisherman? (\textbf{\ta{மீனவர்}}, \emph{mīṉavar}.) A potter? (\textbf{\ta{குயவர்}}, \emph{kuyavar}.) A barber? (\textbf{\ta{முடிதிருத்துபவர்}}, \emph{muṭitiruttupavar}.) A mason? (\textbf{\ta{கொத்தனார்}}, \emph{kottaṉār}.) A scientist? (\textbf{\ta{விஞ்ஞானி}}, \emph{viññāṉi}.)
\end{tcolorbox}
